% ============================================================
% frontmatter.tex —— 扉页与给读者的一封信
% ============================================================
\begin{titlepage}
\centering
\vspace*{2.5cm}
{\Huge\bfseries 烧掉数学书之后}\\[1.2em]
{\Large 和初中生一起重造数学}\\[2.2em]
{\color{shaohong}\rule{0.62\textwidth}{1.2pt}}\\[1.2em]
{\large 抽象代数\quad ·\quad 代数几何\quad ·\quad 微分几何\quad ·\quad 分析\quad ·\quad 初等数论\quad ·\quad 解析数论}\\[7em]
{\large 〔作者姓名〕\quad 著}\\[1.5em]
{\large 二〇二六年八月}
\vfill
{\small 本书献给所有被"背公式"伤害过、却仍然对世界好奇的人。}
\vspace*{1.5cm}
\end{titlepage}

\chapter*{给读者的一封信}
\addcontentsline{toc}{chapter}{给读者的一封信}

亲爱的读者：

先问你一个问题：加法是谁规定的？

你也许从来没想过这件事。加法就像空气，生下来就在那里，老师教，课本印，考试考，仅此而已。可是请你现在就试一试：自己发明一种运算，随便什么运算。比如，规定 $a \,\heart\, b = a + b + 1$，读作"$a$ 爱上 $b$"。那么 $2 \,\heart\, 3 = 6$，$5 \,\heart\, 5 = 11$。恭喜你，你刚刚发明了一个全世界数学课本上没有（好吧，严格说有，但你的老师多半没讲过）的新运算。它好用吗？它遵守交换律吗？它有没有"什么都不做"的数？这些问题我们第一章以后慢慢玩。

这本书想告诉你一件大事：\textbf{数学不是从天上掉下来砸中数学家的，数学是人一间一间造出来的房子}。负数曾经被骂作"比没有还少的荒谬东西"；分数是因为整数不够用才被迫发明的；连"$0$"都是人类发明"没有"本身之后，才终于能写下 $10$ 这个数。课本把它们写得理所当然，是因为课本懒得讲盖房的过程——而我们这本书，就是要拆掉成品房，带你回到工地上，亲手把每一块砖砌回去。

你会说：我又不是数学家。好消息是：\textbf{在工地上，你和数学家拥有完全一样的工具和许可证}。这条许可证全文如下，请收好：

\begin{center}
\begin{tcolorbox}[colback=shaohong!5!white,colframe=shaohong!85!black,arc=3mm,width=0.92\textwidth]
\centering\bfseries 发明许可证\\[0.5em]
\normalfont 持证人遇到不懂的东西时，有权给它随便起个名字；\\
有权自己规定它怎么算，只要规则不自相矛盾；\\
有权质疑任何课本上的定义，并追问"为什么要这样定"。\\[0.5em]
本许可证永久有效，无需年检。
\end{tcolorbox}
\end{center}

全书六个部分——初等数论、抽象代数、分析、解析数论、微分几何、代数几何——听起来像大学数学系的课程表，吓人。但我承诺：\textbf{你只需要会初中数学}：四则运算、分数、负数、乘方、乘法公式 $(a+b)^2=a^2+2ab+b^2$、解一次二次方程、三角形内角和、直角坐标系、抛物线 $y=x^2$。缺什么工具，我们就在书里现场造：对数在第 19 章发明，三角函数几乎不用。每个吓人的术语出场之前，都会先顶着一个土名字干活——"极限"先叫\Keyword{落脚点}，"导数"先叫\Keyword{瞬间变化率}，"积分"先叫\Keyword{堆积}。等你亲手用过它们，正式名字只是一个顺手换上的体面标签，而不是一堵墙。

每章都有固定的几个栏目：\Keyword{开胃问题}让你先动手玩；\Keyword{发明时刻}记录我们被迫发明了什么；\Keyword{动手实验}里的题目按难度标了星：$\star$ 思考，$\star\star$ 计算，$\bigstar\bigstar\bigstar$ 烧脑（选做）。还有\Keyword{诚实声明}——凡是我偷了懒、铺了地毯、只给结论不给证明的地方，都会明明白白告诉你地毯下面是什么。这本书不装神秘。

你会见到两千年前韩信数兵的方法怎么变成今天保护网购付款的密码锁；会见到刘徽用一把"无限细分"的刀切出圆周率，而这把刀两千年后直接切出了微积分；会见到一只不离开薯片的蚂蚁如何测出世界的弯曲；还会见到一条名叫"椭圆"却和椭圆毫无关系的曲线，如何同时拴住费马大定理和你的手机钱包。

准备好了吗？合上课本，深呼吸。我们烧掉旧地图，是为了画一张自己的。

祝你玩得开心。

\hfill 一个陪你重造数学的人
